\section{La regla de decisión}\label{sec:decision}

El detector (\cod{engine/detector/detector.py}) junta las piezas anteriores y responde, para cada
llamada, una acción. No aprende nada: es un conjunto fijo de reglas con siete parámetros. Esta
sección las da completas, porque es aquí donde se decide qué cuesta un error en cada dirección.

\subsection{Qué pasa dentro de una decisión}

Para cada llamada, en este orden:

\begin{enumerate}
  \item El grafo de archivos calcula $m$, $m_{\mathrm{sd}}$ y \cod{jump} para los recursos de la
    llamada.
  \item El grafo de llamadas agrega la llamada, busca sus aristas de contenido y devuelve sus
    etiquetas $U$ y $S$ y, si corresponde, un camino.
  \item El contador de bits observa la llamada. Si el destino es un destino permitido de una clave,
    lo que el contador sumó para esa clave en esta llamada se deshace.
  \item Se arman las nueve señales, el modelo calcula $\ell_t$ y la suma acumulada pasa a $S_t$.
  \item Se actualiza qué claves están expuestas.
  \item Se aplican las reglas.
  \item Si alguna regla pide rotar, se rota: se le pide el valor nuevo al gestor y se reinician la
    cuenta y la exposición de esa clave.
\end{enumerate}

\subsection{Cuándo una clave está expuesta}

Una clave está \emph{expuesta} si la sesión pudo haberla leído: una llamada tocó su archivo, o el
resultado de una llamada contiene 8 caracteres seguidos de su valor vigente. Las tres causas de
rotación exigen que la clave esté expuesta. Tiene sentido: no se rota una clave que nadie leyó, por
mucho que una llamada se parezca a un envío.

Rotar una clave le quita la exposición. El valor nuevo no lo ha leído nadie, y hasta que la sesión no
vuelva a leerlo no hay nada que proteger.

\subsection{Las reglas}

La \cref{fig:decision} las resume y el \cref{alg:decide} las da con precisión.

\begin{figure}[ht]
\centering
% La regla de decisión, en el orden en que se evalúa.
\begin{tikzpicture}[font=\footnotesize\sffamily,
    preg/.style={caja,minimum width=7.9cm,text width=7.5cm,align=left,inner sep=5pt},
    res/.style={caja,minimum width=4.0cm,text width=3.7cm,align=left,inner sep=5pt}]
  \node[preg] (r1) at (0,0)
    {\textbf{1.} ¿La llamada saca datos, subieron los bits \emph{confirmados} de una clave y el destino
     no es uno de los suyos?};
  \node[preg,below=5mm of r1] (r2)
    {\textbf{2.} Para cada clave \emph{expuesta}: ¿salió una fracción $\geq \lrot$, o la evidencia es
     $S \geq \brot$, o pasaron $T$ llamadas desde su última rotación?};
  \node[preg,below=5mm of r2] (r3)
    {\textbf{3.} ¿La evidencia es $S > b$ y hay un camino de contenido: no confiable, clave leída,
     sumidero?};
  \node[preg,below=5mm of r3] (r4)
    {\textbf{4.} ¿Alguna clave que hay que rotar ya rotó 3 veces en las últimas 200 llamadas?};
  \node[caja,below=5mm of r4,minimum width=7.9cm] (ok) {Si nada de lo anterior: \textbf{permitir}};

  \node[res,draw=tinta,fill=tinta,text=white,right=9mm of r1] (a1) {\textbf{Bloquear} y rotar esa clave};
  \node[res,cajaoro,text width=3.7cm,align=left,right=9mm of r2] (a2) {\textbf{Rotar} la clave y dejar pasar la llamada};
  \node[res,draw=tinta,fill=tinta,text=white,right=9mm of r3] (a3) {\textbf{Bloquear}};
  \node[res,cajagris,text width=3.7cm,align=left,right=9mm of r4] (a4) {No rotar más: \textbf{modo escalado}, se bloquean sus salidas};

  \foreach \a/\b in {r1/a1,r2/a2,r3/a3,r4/a4} {\draw[flecha] (\a.east) -- node[above,font=\scriptsize] {sí} (\b.west);}
  \draw[flecha] (r1) -- node[right,font=\scriptsize] {no} (r2);
  \draw[flecha] (r2) -- (r3);
  \draw[flecha] (r3) -- (r4);
  \draw[flecha] (r4) -- (ok);
\end{tikzpicture}

\caption{La regla de decisión. Las reglas 2 y 3 se evalúan solo si la regla 1 no bloqueó. Si se
cumplen a la vez una rotación y un bloqueo, la acción es \cod{rotate+block}.}
\label{fig:decision}
\end{figure}

\begin{description}
  \item[Regla 1: bloqueo determinista.] La llamada es un sumidero, los bits \emph{confirmados} de una
    clave subieron en esta llamada y el destino no es uno de los permitidos de esa clave (o no se
    pudo identificar). Se bloquea la llamada y se rota la clave. Es la única regla que no exige
    exposición ni evidencia: un tramo de 8 o más caracteres seguidos de una clave, saliendo hacia
    donde no debe, no admite otra lectura. Las candidatas nunca disparan esta regla.
  \item[Regla 2: rotar.] Para cada clave expuesta, la primera de estas condiciones que se cumpla:
    \begin{itemize}
      \item \emph{por el contador:} la fracción filtrada, sumando confirmados y candidatos, es
        $\geq \lrot$;
      \item \emph{por la evidencia:} $S_t \geq \brot$;
      \item \emph{por el período:} pasaron $T$ llamadas desde la última rotación de esa clave.
    \end{itemize}
    La clave rota y la llamada \emph{se deja pasar}. Si la regla la dispara la evidencia, rotan todas
    las claves expuestas en ese momento, porque la evidencia es de la sesión y no de una clave.
  \item[Regla 3: bloqueo por evidencia.] $S_t > b$ y el grafo de llamadas entregó un camino de
    contenido. Se bloquea.
  \item[Regla 4.] Si nada de lo anterior se cumple, se permite.
\end{description}

\begin{algorithm}[ht]
\caption{La regla de decisión (\cod{Detector.decide}), después de medir.}\label{alg:decide}
\KwIn{la llamada $c$, su destino $d$, la evidencia $S$, el camino $\pi$ del grafo de llamadas; por cada
clave $s$: bits confirmados antes y después, fracción $\lambda(s)$, exposición, llamadas $e(s)$ desde su
última rotación}
$\mathit{rotar} \leftarrow \emptyset$; $\mathit{bloquear} \leftarrow$ falso\;
\If(\Com*[f]{regla 1}){$c$ saca datos}{
  \ForEach{clave $s$ cuyos bits confirmados subieron y $d \notin \mathit{permitidos}(s)$}{
    $\mathit{bloquear} \leftarrow$ verdadero; agregar $s$ a $\mathit{rotar}$\;
  }
}
\If{no $\mathit{bloquear}$}{
  \ForEach(\Com*[f]{regla 2}){clave $s$ expuesta}{
    \lIf{$\lambda(s) \geq \lrot$ \textbf{o} $S \geq \brot$ \textbf{o} $e(s) \geq T$}{agregar $s$ a $\mathit{rotar}$}
  }
  \lIf(\Com*[f]{regla 3}){$S > b$ \textbf{y} $\pi \neq \emptyset$}{$\mathit{bloquear} \leftarrow$ verdadero}
}
\ForEach(\Com*[f]{tope de rotaciones}){$s \in \mathit{rotar}$}{
  \lIf{$s$ ya rotó $3$ veces en las últimas $200$ llamadas}{quitar $s$ de $\mathit{rotar}$; $s$ pasa a modo escalado}
}
\If{$c$ saca datos \textbf{y} hay una clave $s$ en modo escalado con $d \notin \mathit{permitidos}(s)$}{
  $\mathit{bloquear} \leftarrow$ verdadero\;
}
\Return \cod{rotate+} si $\mathit{rotar} \neq \emptyset$, seguido de \cod{block} si $\mathit{bloquear}$ y de \cod{allow} si no\;
\end{algorithm}

\subsection{El tope de rotaciones y el modo escalado}\label{sec:escalado}

Rotar no puede ser la respuesta indefinida. Un atacante que relee la clave después de cada rotación
obligaría a rotar sin fin, y cada rotación le deja sacar un poco de la versión nueva. Por eso una clave
puede rotar a lo más 3 veces dentro de una ventana de 200 llamadas.

Al cuarto pedido la clave no rota: entra en \emph{modo escalado}. Desde ahí se bloquea todo sumidero
hacia un destino que no sea uno de sus permitidos, con el motivo «demasiadas rotaciones: se bloquean
las salidas y se avisa a una persona». La clave sale del modo escalado cuando pasan 200 llamadas sin
que ninguna regla pida rotarla.

Es el único momento en que Alicanto deja de preferir la rotación al bloqueo, y lo hace a propósito:
contra un adversario persistente, termina deteniendo las salidas. El aviso a una persona es hoy el
texto del motivo; no hay un canal de notificación implementado. El costo de este modo sobre el trabajo
normal no está medido.

\subsection{La invariante que acota el período}\label{sec:invariante}

El período $T$ no se puede elegir libremente. Hay un adversario que conoce las reglas del contador y
juega contra ellas:

\begin{itemize}
  \item envía pedazos de $k - 1 = 7$ caracteres, el más largo que no se confirma;
  \item espera $W + 1 = 51$ llamadas entre un pedazo y el siguiente, para que cada candidata caduque
    antes de que llegue la próxima.
\end{itemize}

Así el contador nunca ve más de un pedazo a la vez: 28 bits, bajo el umbral del \qty{30}{\percent} de
la demo. Lo único que lo detiene es la rotación periódica. Entre dos rotaciones periódicas alcanza a
hacer $\lceil T / (W+1) \rceil$ envíos, de modo que la configuración solo es segura si
\begin{equation}
  \left\lceil \frac{T}{W + 1} \right\rceil \cdot (k - 1) \cdot \log_2 |A_s| \;\leq\; H - \beta .
  \label{eq:invariante}
\end{equation}
Con claves hexadecimales de 160 bits, el lado izquierdo vale 112 con $T = 200$, 84 con $T = 150$ y 56
con $T = 75$, contra un límite de 96. El valor inicial del diseño era 200 y no cumplía: ese adversario
sacaba 112 bits. Se bajó a 150 durante la construcción, y la condición quedó escrita en el código
(\cod{Params.check}): un detector con parámetros que la violan no arranca, y la búsqueda de parámetros
descarta esas combinaciones antes de evaluarlas.

\begin{limite}
\textbf{La invariante no cubre los pedazos de menos de 4 caracteres.} El contador no los cuenta, así
que para ellos el argumento no aplica. Sin el modelo de evidencia, un carácter por llamada saca la
clave entera antes de la primera rotación periódica. La \cref{sec:fragmentos} lo mide.
\end{limite}

Con la configuración evaluada la invariante importa menos: como $\lrot = \num{0.1}$ equivale a 16 bits,
cualquier pedazo que el contador cuente (4 caracteres o más) ya dispara la rotación.

\subsection{Los parámetros}\label{sec:parametros}

\begin{table}[ht]
\caption{Los parámetros de la regla, en las dos configuraciones. El \apref{ap:parametros} agrega las
constantes del contador y de los grafos.}\label{tab:params}
\small
\begin{tabularx}{\textwidth}{@{}P{2.6cm}R{1.6cm}R{1.9cm}L@{}}
\toprule
Parámetro & Demo & Evaluada & Qué controla \\
\midrule
$\lrot$ (\cod{lambda_rot}) & \num{0.3} & \num{0.1} & Fracción filtrada que dispara la rotación: 48 o
  16 de 160 bits \\
$\brot$ (\cod{b_rot}) & \num{4.0} & \num{2} & Evidencia acumulada que dispara la rotación \\
$b$ & \num{9.2} & \num{4} & Evidencia acumulada que, con un camino de contenido, bloquea \\
$T$ & 150 & 75 & Llamadas entre rotaciones periódicas de una clave expuesta \\
$\beta$ (\cod{beta_bits}) & 64 & 64 & Seguridad residual exigida \\
\cod{max_rotations} & 3 & 3 & Rotaciones de una misma clave dentro de la ventana \\
\cod{rotation_window} & 200 & 200 & Largo de la ventana, en llamadas \\
Modelo de evidencia & no & sí & Sin modelo, $\ell_t = 0$ y $S_t = 0$ siempre \\
\bottomrule
\end{tabularx}
\end{table}

La configuración de la demo no carga el modelo de evidencia. Con $S_t = 0$ siempre, la condición por
evidencia de la regla 2 y la regla 3 no se disparan nunca: quedan el contador, el período y el bloqueo
determinista. Es lo que se quiere mostrar en una demo (el goteo sube el medidor hasta que rota), y es
también la configuración del \esp{hook} de Claude Code.

Con la configuración evaluada y el modelo saturado de la \cref{sec:evidencia}, los umbrales se leen
así: como $\ell$ vale $+3$ o $-3$, $\brot = 2$ se cruza con \emph{una sola} llamada marcada, y $b = 4$
exige dos seguidas. La \cref{sec:congelamiento} cuenta cómo se eligieron y qué tan determinados
quedaron por los datos.

\subsection{Modos de medición}

El mismo detector corre en tres variantes que usa la evaluación.

\begin{description}
  \item[Solo medición] (\cod{enforce=False}). Mide todo igual, pero la acción es siempre
    \cod{allow} y no rota. Es el panel «sin protección» de la demo, y el modo con que se extraen las
    señales para ajustar el modelo.
  \item[Sin rotación] (\cod{rotation_enabled=False}). Ninguna regla rota; los bloqueos siguen. Es la
    ablación que mide cuánto aporta rotar.
  \item[Sin evidencia] (\cod{model=None}). Los mismos parámetros, sin modelo. Es la ablación que mide
    cuánto aporta la evidencia.
\end{description}

\subsection{Lo que dice cada decisión}

Cada decisión lleva lo necesario para explicarla sin abrir el código:

\begin{itemize}
  \item \cod{reason}: una frase con la regla y sus números, por ejemplo «Rotación de SECRET\_1:
    fracción filtrada 0.30 $\geq$ 0.3» o «Rotación de SECRET\_1: expuesto y evidencia S=3.00 $\geq$ 2».
  \item \cod{leak}: por clave, los bits y la fracción que habían salido \emph{antes} de rotar.
  \item \cod{path}: las llamadas que explican la decisión, en orden: la primera lectura de la clave,
    los envíos que aportaron bits y la llamada actual.
  \item \cod{S} y \cod{latency_ms}: la evidencia acumulada y lo que tardó decidir.
\end{itemize}
